\clearpage
\SourceLeaf{39}

\begin{center}
\textbf{[MECHANISCHER SATZ ÜBER DIE WURZELN EINER GANZEN
FUNKTION $f(x)$ UND IHRER ABLEITUNG $f'(x)$.]}
\end{center}

\noindent\hspace*{\parindent}Sind $\begin{array}{c}a,\,b,\,c,\,.\,.,\,m,\,n\\ a',\,b',\,c',\,.\,.,\,m'\end{array}$ die
Wurzeln der Gleichung $\begin{array}{c}f\,x=0\\ f'\,x=0,\end{array}$ wobei $f'x=\dfrac{d\,f\,x}{d\,x}$,
und bezeichnen dieselben Buchstaben die entsprechenden Punkte in der Ebene, so herrscht, wenn
man sich in $a,\,b,\,c,\,.\,.,\,m,\,n$ gleiche abstoßende oder anziehende Massen denkt, die umgekehrt proportional
zur Entfernung wirken, in $a',\,b',\,c',\,.\,.,\,m'$ Gleichgewicht.

\Needspace{9\baselineskip}
\begin{center}
\textbf{BEMERKUNGEN ZU DEN BEIDEN VORSTEHENDEN ALGEBRAISCHEN NOTIZEN.}
\end{center}

\noindent\hspace*{\parindent}Die erste der beiden hier mitgeteilten Notizen stammt wahrscheinlich aus den Jahren 1826--28. Sie
dient dazu, ein System von fünf linearen Gleichungen mit einer besonders einfachen, symmetrisch
aufgebauten Determinante mithilfe der fünften Einheitswurzeln zu lösen. Beim Beweis der
sich ergebenden Darstellungen für die fünf Größen $p,\,q,\,r,\,s,\,t$ beachte man,
dass die Norm der komplexen Zahl $a+b\varepsilon+c\varepsilon\varepsilon+d\varepsilon^3+e\varepsilon^4$
sich in der Form $\frac{1}{4}(5M-m\mathfrak{M})$ darstellen lässt. 

Die zweite Notiz hat \textsc{Gauss} wahrscheinlich im Jahr 1846 aufgezeichnet. Zum Beweis
richte man die Ebene so aus, dass eine der Lösungen von $f'(x)=0$ im Ursprung liegt.
Dann gilt:

\[
\begin{aligned}
\frac{1}{a}+\frac{1}{b}+\frac{1}{c}+\cdots+\frac{1}{n}&=0,\\[3pt]
\frac{a_0}{a\,a_0}+\frac{b_0}{b\,b_0}+\frac{c_0}{c\,c_0}+\cdots+\frac{n_0}{n\,n_0}&=0,
\end{aligned}
\]

\noindent wobei $a_0$ zu $a$, $b_0$ zu $b$ usw. komplex konjugiert sein sollen. Die mechanische Deutung der
letzten Gleichung liefert den Satz von \textsc{Gauss}.
\par\nobreak\hfill\textsc{Fricke.}\par

